\chapter{Das Hauptfenster}\label{kap:hauptfenster}
\index{Hauptfenster}

Dieses Kapitel stellt alle Bereiche des Hauptfensters vor und beschreibt jeden Menüeintrag.
Wie man sich in den Ansichten bewegt und Objekte auswählt, steht in
\autoref{kap:navigation}; Zeichnen und Modellieren in \autoref{kap:modellieren}.

\section{Überblick}\label{sec:fenster-ueberblick}

\abbildung[\linewidth]{main-window-shaded.png}{Das Hauptfenster mit dem Beispielmodell
\datei{demo-parts.step} in der Darstellung \gui{Shaded}. Links die Werkzeugspalte, in der Mitte
die vier Ansichten, oben die Werkzeugleiste, rechts die Docks \gui{Properties} und
\gui{Model}, unten Befehlszeile und Statusleiste.}{fig:hauptfenster}

Das Hauptfenster (\autoref{fig:hauptfenster}) besteht aus diesen Bereichen:

\begin{referenz}{>{\raggedright\arraybackslash}p{3.6cm}>{\raggedright\arraybackslash}p{2.4cm}X}
\toprule
\textbf{Bereich} & \textbf{Lage} & \textbf{Zweck} \\
\midrule
\endhead
Menüleiste & oben & \menu{File}, \menu{Edit}, \menu{View}, \menu{Plugins}, \menu{Help}
  (\autoref{sec:menues}); auf dem Mac in der Menüleiste am oberen Bildschirmrand. \\
Werkzeugleiste \gui{View} & oben & Darstellungsmodi, Zoom, Ein- und Ausblenden
  (\autoref{sec:werkzeugleiste}). \\
Werkzeugspalte \gui{Tools} & links & alle Befehle in aufklappbaren Gruppen, dazu je eine
  Gruppe pro Plugin (\autoref{sec:werkzeugspalte}). \\
Vier Ansichten & Mitte & Front, Top, Left und Perspective (\autoref{sec:ansichten}). \\
Docks & rechts, unten oder frei & \gui{Properties}, \gui{Layers}, \gui{Model},
  \gui{Results}, \gui{Plugin}, \gui{Plot}, \gui{Table}, \gui{Image}, \gui{Plugin log},
  \gui{Node workflow} (\autoref{sec:docks}). \\
Befehlszeile & unten & Befehle und Koordinaten eintippen (\autoref{sec:befehlszeile}). \\
Statusleiste & ganz unten & Meldungen, Cursorkoordinaten, aktive Ansicht, Fortschritt von
  Rechnungen, \gui{Ortho} und Fangschalter (\autoref{sec:statusleiste}). \\
\bottomrule
\end{referenz}

\paragraph{Fenstertitel.} Der Titel nennt den Dateinamen des Dokuments (\enquote{Untitled}
für ein neues) und den Programmnamen, z.\,B.\ \enquote{optik.rttf -- Raytatouille Fusion}.
Hat das Dokument ungespeicherte Änderungen, erscheint ein Sternchen hinter dem Namen (auf dem
Mac ein Punkt im roten Schließknopf).

\paragraph{Anordnung merken.} Größe und Lage des Fensters sowie die Anordnung aller Docks
speichert das Programm beim Beenden und stellt sie beim nächsten Start wieder her
(\autoref{sec:einstellungen}).

\section{Das dunkle Design}\label{sec:design}
\index{dunkles Design}

\produkt{} verwendet unabhängig von der Einstellung des Betriebssystems ein dunkles Design:
dunkelgraue Flächen, helle Schrift und einen blauen Hervorhebungston. Unter macOS ist das die
native dunkle Darstellung des Systems, unter Windows und Linux der Qt-Stil \enquote{Fusion}
mit einer eigenen dunklen Farbpalette, damit auch Menüs, Dialoge und Beschriftungen gut
lesbar sind. Die Ansichten haben einen dunklen Hintergrund; die perspektivische Ansicht ist
etwas heller als die drei Parallelansichten.

Die Farben im Modell haben eine feste Bedeutung:

\begin{referenz}{>{\raggedright\arraybackslash}p{4.2cm}X}
\toprule
\textbf{Farbe} & \textbf{Bedeutung} \\
\midrule
\endhead
Ebenenfarbe (Standard hellgrau) & normale Objekte; die Farbe kommt von der Ebene oder vom
  Objekt selbst (\autoref{sec:ebenen}). \\
Gelb & ausgewählte Objekte und Teilobjekte. \\
Cyan & das Objekt bzw.\ Teilobjekt unter dem Mauszeiger (Vorschau der Auswahl). \\
Orange & Gummiband: Vorschau dessen, was ein laufender Befehl erzeugen wird. \\
Grün & Fangsymbol am Cursor (\autoref{kap:modellieren}). \\
Rot, Grün, Blau & die Achsen X, Y und Z -- im Achsenkreuz unten links jeder Ansicht und
  als Achsen der Konstruktionsebene im Raster. \\
Blauer Rahmen, durchgezogen / grüner Rahmen, gestrichelt & Auswahlrechteck als Fenster bzw.\
  als Kreuzen (\autoref{sec:rechteckauswahl}). \\
\bottomrule
\end{referenz}

\section{Die Werkzeugspalte}\label{sec:werkzeugspalte}
\index{Werkzeugspalte}

Die Werkzeugspalte am linken Rand (Dock \gui{Tools}) enthält alle Zeichen-, Modellier- und
Analysebefehle als Knöpfe, geordnet in Gruppen. Sie nimmt beim ersten Start etwa ein Zehntel
der Fensterbreite ein (mindestens 96 Pixel); ihre Breite lässt sich am Rand zu den Ansichten
mit der Maus ändern. Die Knöpfe ordnen sich dann in so vielen Spalten an, wie hineinpassen.

\paragraph{Gruppen ein- und ausklappen.} Jede Gruppe hat eine Kopfzeile mit Pfeil und
Namen. Ein Klick darauf klappt die Gruppe zu ($\triangleright$) oder auf
($\triangledown$). Es können beliebig viele Gruppen gleichzeitig offen sein; welche offen
sind, merkt sich das Programm. Anfangs sind \gui{Draw} und \gui{Curves} sowie die Gruppen der
Plugins aufgeklappt, alle übrigen zugeklappt.

\paragraph{Tooltips.} Hält man den Mauszeiger über einen Knopf, erscheint ein Tooltip mit dem
vollen Namen des Befehls und darunter dem Alias, den Sie tippen können, z.\,B.\
\enquote{Line} und \enquote{Type: l + Enter}; die Statusleiste zeigt den Namen ebenfalls. Die
Knöpfe tragen derzeit kurze Textkürzel statt Symbolen (\gui{Pt} für Point, \gui{Ln} für Line
\dots).

Der Alias ist kein Tastenkürzel: Tippen Sie in einer Ansicht \code{l} und drücken Sie
\keys{Enter} -- der Buchstabe landet in der Befehlszeile, erst \keys{Enter} startet den Befehl.
Jeder Befehl hat einen solchen Alias (\autoref{kap:befehle}). Ein echtes Tastenkürzel hat nur
\gui{Delete} (\keys{Entf}); sein Tooltip nennt es in Klammern.

\paragraph{Befehl starten.} Ein Klick auf einen Knopf startet den Befehl. Die Befehlszeile
zeigt dann, welche Eingabe erwartet wird (\autoref{sec:befehlszeile}). Ein laufender Befehl wird
durch einen neuen ersetzt.

Die festen Gruppen und ihre Befehle (Einzelheiten zu jedem Befehl in
\autoref{kap:modellieren} und \autoref{kap:befehle}):

\begin{referenz}{>{\raggedright\arraybackslash}p{2.2cm}X}
\toprule
\textbf{Gruppe} & \textbf{Befehle (Knopfbeschriftung: Name; in Klammern, was der Tooltip zu
  tippen empfiehlt)} \\
\midrule
\endhead
\gui{Draw} & Pt: Point (\code{pt}), Ln: Line (\code{l}), PL: Polyline (\code{pl}), Rec:
  Rectangle (\code{rec}), Pol: Polygon (\code{pol}) \\
\gui{Curves} & Ci: Circle (centre, radius) (\code{c}), C3: Circle (3 points) (\code{c3}), Arc:
  Arc (\code{a}), Hx: Helix (\code{h}), Sp: Interpolated spline (\code{sp}), Off: Offset curve
  (\code{offset}), Fil: Fillet curves (\code{fillet}) \\
\gui{Surfaces} & Pln: Plane (\code{plane}), Ext: Extrude (\code{ext}), Rev: Revolve
  (\code{rev}), Lft: Loft (\code{loft}), Sw1: Sweep 1 rail (\code{sweep1}), Sw2: Sweep 2 rails
  (\code{sweep2}), Cns: Coons patch (\code{coons}), Gor: Gordon surface (\code{gordon}), OfS:
  Offset surface (\code{offsetsrf}), FiS: Fillet surfaces (\code{filletsrf}), ChS: Chamfer surfaces
  (\code{chamfersrf}), Cap: Cap planar
  holes (\code{cap}) \\
\gui{Solids} & Box (\code{b}), Sph: Sphere (\code{sph}), Cyl: Cylinder (\code{cyl}), Con: Cone
  (\code{cone}), Tor: Torus (\code{tor}), $\cup$: Boolean union (\code{bu}), $-$: Boolean
  difference (\code{bd}), $\cap$: Boolean intersection (\code{bi}) \\
\gui{Transform} & Mv: Move (\code{m}), Cp: Copy (\code{co}), Rot: Rotate (\code{ro}), Scl:
  Scale (\code{sc}), Mir: Mirror (\code{mi}), ArL: Linear array (\code{array}), ArG: Grid
  array (\code{array\_grid}), ArP: Polar array (\code{arraypolar}) \\
\gui{Edit} & Trm: Trim (\code{tr}), Spl: Split (\code{split}), Jn: Join (\code{j}), Exp:
  Explode (\code{x}), Del: Delete (Tastenkürzel \keys{Entf}) \\
\gui{Analyze} & Dst: Distance (\code{di}), Ang: Angle (\code{ang}), Crv: Curvature
  (\code{curv}), Are: Area (\code{area}), Vol: Volume / mass properties (\code{volume}), Int:
  Intersection check (\code{intersect}), ACe: Area centroid, GCo: Continuity (\code{gcon}), Zeb:
  Zebra, EMp: Environment map, CAn: Curvature analysis, Dft: Draft angle analysis, AOf: Analysis
  off, Chk: Check, SBO: Select bad objects \\
\bottomrule
\end{referenz}

\paragraph{Gruppen der Plugins.} Unter den festen Gruppen folgt für jedes geladene Plugin eine
Gruppe mit seinem Namen, deren Knöpfe die Aktionen des Plugins starten. Die Knöpfe tragen die
ersten drei Buchstaben der Aktion; der Tooltip nennt den vollen Namen. Dieselben Aktionen
stehen im Untermenü des Plugins im Menü \menu{Plugins} und als Knöpfe im Dock \gui{Plugin}.
Mit \menu{Plugins,Show in the side panels} lässt sich die Gruppe eines Plugins ausblenden
(\autoref{sec:menu-plugins}).

\paragraph{Spalte schließen, verschieben, zurückholen.} Die Werkzeugspalte ist ein Dock: Sie
lässt sich an ihrer Titelleiste an den rechten Fensterrand ziehen oder mit dem Schließknopf
ausblenden. \menu{View,Tools} holt sie zurück.

\section{Die vier Ansichten}\label{sec:ansichten}
\index{Ansicht}

Die Mitte des Fensters teilen sich vier Ansichten im Raster $2 \times 2$:

\begin{referenz}{>{\raggedright\arraybackslash}p{2.4cm}>{\raggedright\arraybackslash}p{1.8cm}>{\raggedright\arraybackslash}p{3.3cm}X}
\toprule
\textbf{Ansicht} & \textbf{Lage} & \textbf{Blickrichtung} & \textbf{Bildschirm rechts / oben;
  Projektion} \\
\midrule
\endhead
\gui{Front (XZ)} & oben links & entlang $+Y$ & $X$ / $Z$; parallel \\
\gui{Top (XY)} & oben rechts & entlang $-Z$ (von oben) & $X$ / $Y$; parallel \\
\gui{Left (YZ)} & unten links & entlang $+X$ (von der $-X$-Seite) & $-Y$ / $Z$; parallel \\
\gui{Perspective} & unten rechts & schräg von vorn rechts oben, aus Richtung $(+X, -Y, +Z)$ &
  $Z$ zeigt nach oben; perspektivisch \\
\bottomrule
\end{referenz}

\paragraph{Kopfzeile.} Oben links in jeder Ansicht steht ihr Titel, z.\,B.\ \enquote{Top
(XY) $\cdot$ Shaded $\blacktriangledown$}: der Name der Ansicht, der Darstellungsmodus und ein
Pfeil. Weicht die Projektion vom Standard der Ansicht ab, steht dahinter \enquote{(parallel)}
bzw.\ \enquote{(perspective)}. Ein Klick auf die Kopfzeile öffnet das Ansichtsmenü
(\autoref{sec:ansichtsmenue}): andere Standardansicht, Projektion umschalten, Zoom und
Maximieren.

\paragraph{Achsenkreuz und Raster.} Unten links zeigt jede Ansicht ein Achsenkreuz mit den
Achsen X (rot), Y (grün) und Z (blau). Im Hintergrund liegt das Raster der
Konstruktionsebene, dessen Linienabstand sich dem Zoom anpasst
(\autoref{sec:konstruktionsebene}).

\paragraph{Aktive Ansicht.} Die zuletzt angeklickte Ansicht ist \emph{aktiv}; sie hat einen
hervorgehobenen Rahmen, und die Statusleiste nennt sie (\enquote{Active view: Perspective}).
Darstellungsmodus, Zoom und Maximieren aus Werkzeugleiste und Menü wirken auf die aktive
Ansicht. Beim Start ist die perspektivische Ansicht aktiv.

\paragraph{Größe der Ansichten.} Die Trennlinien zwischen den Ansichten lassen sich mit der
Maus verschieben; die senkrechte Trennlinie bleibt dabei in beiden Reihen gleich.
\keys{Leertaste} vergrößert die aktive Ansicht auf die ganze Fläche und stellt beim
nächsten Druck das Raster wieder her (\autoref{sec:maximieren}).

\paragraph{Farblegende.} Zeigen Ergebnisse Werte farbig an -- ein Feld auf einer Fläche oder
nach einer Größe gefärbte Strahlen --, erscheint in den Ansichten eine Farblegende mit dem
Namen des Ergebnisses, der Farbskala, dem Wertebereich und der Einheit. Sie gilt für das
erste sichtbare Ergebnis dieser Art.

\subsection{Darstellungsmodi}\label{sec:darstellungsmodi}
\index{Darstellungsmodus}

Jede Ansicht hat ihren eigenen Darstellungsmodus:

\begin{referenz}{>{\raggedright\arraybackslash}p{2.2cm}X}
\toprule
\textbf{Modus} & \textbf{Darstellung} \\
\midrule
\endhead
\gui{Wireframe} & Drahtmodell: Kurven, Kanten und Isolinien der Flächen, keine gefüllten
  Flächen. Alles ist sichtbar, auch was hinten liegt. Flächen werden über ihre Isolinien
  ausgewählt. \\
\gui{Surfaces} & Flächen in ihrer Farbe ohne Beleuchtung, mit Kanten. Verdeckte Kurven und
  Kanten sind nicht sichtbar und nicht wählbar. \\
\gui{Shaded} & Beleuchtete Flächen (zwei Lichter am Betrachter, Rückseiten dunkler) mit
  Kanten. Am besten für die räumliche Wirkung. \\
\bottomrule
\end{referenz}

\abbildung[0.9\linewidth]{main-window-wireframe.png}{Dasselbe Modell wie in
\autoref{fig:hauptfenster} in der Darstellung \gui{Wireframe}: Kanten und Isolinien, keine
Flächen.}{fig:wireframe}

Umschalten: Klick auf \gui{Wireframe}, \gui{Surfaces} oder \gui{Shaded} in der Werkzeugleiste
oder im Menü \menu{View}. Das wirkt auf die aktive Ansicht; mit gedrückter \keys{Shift}-Taste
geklickt auf alle vier Ansichten. Der Modus der aktiven Ansicht ist in der Werkzeugleiste
eingedrückt. Beim Start sind alle Ansichten im Modus \gui{Shaded}; mit der Option
\code{-{}-display-mode} lässt sich ein anderer Startmodus wählen (\autoref{sec:kommandozeile}).

\begin{tipp}
Der Darstellungsmodus beeinflusst auch das Auswählen und den Fang: In \gui{Surfaces} und
\gui{Shaded} verdecken Flächen, was hinter ihnen liegt. Um eine verdeckte Kante zu treffen,
schalten Sie die Ansicht auf \gui{Wireframe}.
\end{tipp}

\section{Die Werkzeugleiste}\label{sec:werkzeugleiste}
\index{Werkzeugleiste}

Die Werkzeugleiste \gui{View} unter der Menüleiste ist fest und enthält:

\begin{referenz}{>{\raggedright\arraybackslash}p{3.0cm}>{\raggedright\arraybackslash}p{2.6cm}X}
\toprule
\textbf{Knopf} & \textbf{Taste} & \textbf{Wirkung} \\
\midrule
\endhead
\gui{Wireframe} & -- & Drahtmodell in der aktiven Ansicht; \keys{Shift}+Klick: alle
  Ansichten. \\
\gui{Surfaces} & -- & flache Flächen mit Kanten; \keys{Shift}+Klick: alle Ansichten. \\
\gui{Shaded} & -- & beleuchtete Flächen mit Kanten; \keys{Shift}+Klick: alle Ansichten. \\
\gui{Zoom extents} & \keys{Z}, dann \keys{E} & Das ganze Modell füllt die aktive Ansicht;
  \keys{Shift}+Klick: alle Ansichten (\autoref{sec:zoom-extents}). \\
\gui{Zoom selected} & \keys{Z}, dann \keys{S} & Die ausgewählten Objekte, Flächen, Kanten oder
  Punkte füllen die aktive Ansicht; ohne Auswahl wie \gui{Zoom extents}. \keys{Shift}+Klick:
  alle Ansichten. \\
\gui{Hide} & \keys{H} & blendet die ausgewählten Objekte aus und hebt ihre Auswahl auf
  (\autoref{sec:ausblenden}). \\
\gui{Unhide} & \keys{Shift+H} & zeigt alle mit \gui{Hide} oder \gui{Isolate} ausgeblendeten
  Objekte wieder. \\
\bottomrule
\end{referenz}

Die Tooltips der Knöpfe nennen Wirkung und Tastenkürzel.

\section{Die Docks}\label{sec:docks}
\index{Dock}

Docks sind Teilfenster, die am Rand des Hauptfensters andocken. Jedes Dock lässt sich

\begin{itemize}
  \item an seiner Titelleiste an einen anderen Rand ziehen,
  \item auf ein anderes Dock ziehen, sodass beide als Reiter übereinander liegen (unten im
    Dock erscheinen dann Reiter zum Umschalten),
  \item mit dem Knopf neben dem Schließknopf vom Fenster lösen (\enquote{schwebend}) und mit
    einem Doppelklick auf die Titelleiste wieder andocken,
  \item mit dem Schließknopf schließen und über das Menü \menu{View} wieder öffnen.
\end{itemize}

Anfangs sind die Docks so angeordnet: rechts oben \gui{Properties}, darunter als Reiter
\gui{Layers}, \gui{Model}, \gui{Results} und \gui{Plugin} (\gui{Model} vorn). \gui{Plot},
\gui{Table}, \gui{Plugin log} und \gui{Node workflow} liegen unten und sind zunächst
verborgen; \gui{Image} ist ein schwebendes Fenster, ebenfalls zunächst verborgen. Diese Docks
erscheinen von selbst, sobald sie etwas zu zeigen haben.

\subsection{Properties}\label{sec:dock-properties}

Zeigt die Eigenschaften der Auswahl: bei einem einzelnen Objekt Ebene, Id, Typ und die
geometrischen Daten (Grad, Kontrollpunkte, Knoten, Länge, Fläche, Volumen \dots) sowie seine
Attribute; bei mehreren Objekten, wie viele welcher Art ausgewählt sind. Mit \gui{Add
attribute...} werden freie Attribute angelegt. Ausführlich in \autoref{sec:eigenschaften}.

\subsection{Layers}

Die Ebenen des Dokuments mit Sichtbarkeit, Name und Farbe; \gui{New layer} legt eine Ebene
an, \gui{Move selection here} verschiebt die Auswahl auf die markierte Ebene. Ausführlich in
\autoref{sec:ebenen}.

\subsection{Model}

Der Modellbaum: Gruppen und Objekte mit Name und Typ, mit Kontrollkästchen für die
Sichtbarkeit. Ein Klick wählt aus (in beide Richtungen: Auswahl in der Ansicht markiert die
Zeile), das Kontextmenü gruppiert und blendet aus. Ausführlich in \autoref{sec:modellbaum}.

\subsection{Results}\label{sec:dock-results}
\index{Ergebnis}

Hier stehen die Ergebnisse der Plugin-Rechnungen. Jede Rechnung erzeugt einen \emph{Lauf} (eine
Zeile oberster Ebene, in der Spalte \gui{Type} steht, woher er kommt); darunter die einzelnen
Ergebnisse mit ihrer Art:

\begin{referenz}{>{\raggedright\arraybackslash}p{3.0cm}X}
\toprule
\textbf{Art (\gui{Type})} & \textbf{Anzeige} \\
\midrule
\endhead
\gui{Rays / lines} & Linien in den Ansichten, z.\,B.\ Strahlen; der Tooltip nennt die Zahl
  der Polylinien. \\
\gui{Field on face} & farbige Werte auf einer Fläche (Bestrahlungsstärke, Feld), mit
  Farblegende. \\
\gui{Curve} & Kurve im Dock \gui{Plot}. \\
\gui{Bars} & Balkendiagramm im Dock \gui{Plot}. \\
\gui{Image} & Bild im Fenster \gui{Image}. \\
\gui{Table} & Tabelle im Dock \gui{Table}. \\
\bottomrule
\end{referenz}

Die Kontrollkästchen schalten einen ganzen Lauf oder einzelne Ergebnisse sichtbar oder
unsichtbar. Das \maus{Kontextmenü} (Rechtsklick)
\begin{itemize}
  \item auf einem Lauf: \gui{Save run…} schreibt alle seine Ergebnisse als PNG- und
    CSV-Dateien in einen Ordner (\autoref{sec:ergebnisse-speichern}); \gui{Delete run}
    entfernt den Lauf mit allen Ergebnissen;
  \item auf einem Ergebnis: \gui{Save…} speichert es als PNG (Bilder, Felder) oder CSV;
  \item auf einem Linien-Ergebnis: \gui{Convert to CAD curves...} fragt
    \enquote{Every n-th ray (of \dots)} -- jeder wievielte Strahl übernommen werden soll --
    und legt die Strahlen als Kurven im Dokument an (ein Rückgängig-Schritt \enquote{Rays to
    curves}). So werden Strahlen zu Geometrie, die gespeichert, exportiert und als
    Fangpunkte benutzt werden kann.
\end{itemize}
Die Schaltfläche \gui{Save…} unter dem Baum speichert die gewählte Zeile (Lauf oder Ergebnis).

Ergebnisse sind keine Geometrie und werden nicht als solche gespeichert. Die neuesten
Ergebnisse der Anzeigeblöcke des Workflows speichert das Programm aber mit dem Dokument; nach
dem Öffnen stehen sie im Lauf \enquote{Saved results} (\autoref{sec:gespeicherte-ergebnisse}).
Ein neues oder geöffnetes Dokument leert das Dock.

\subsection{Plugin}\label{sec:dock-plugin}

Die Bedienseite der Plugins. Oben wählt eine Auswahlliste das Plugin; darunter stehen seine
Beschreibung, seine Parameter in einem Formular und die Knöpfe seiner Aktionen. Parameter
werden mit dem Dokument gespeichert. \gui{Reset parameters} setzt alle Parameter des Plugins
auf die Standardwerte zurück. Während eine Aktion läuft, sind die Aktionsknöpfe gesperrt.
Rot markierte Felder enthalten ungültige Werte; solange es solche gibt, startet keine Aktion.
Läuft ein Plugin nicht, steht statt der Knöpfe \enquote{not running -- see the plugin log}.
Ausführlich in \autoref{kap:plugins}.

\subsection{Plot}\label{sec:dock-plot}

Ein Diagramm für Kurven und Balken aus den Ergebnissen, etwa den Verlauf der Gütefunktion
einer Optimierung, der während der Rechnung wächst. Achsen mit Beschriftung, die Namen der
Kurven als Legende oben links. Das Dock erscheint unten, sobald die erste Kurve vorliegt;
ohne Kurven steht \enquote{No curves (e.g.\ the merit function of an optimisation)} darin.

\subsection{Table}

Tabellen aus den Ergebnissen, z.\,B.\ die Wirkungsgrade der Beugungsordnungen eines Gitters.
Gibt es mehrere, wählt eine Auswahlliste oben die angezeigte; eine neue Tabelle wird
automatisch angezeigt. Die Zellen sind nicht editierbar; Spaltenbreiten lassen sich ziehen.
Das Dock erscheint mit der ersten Tabelle und steht nicht im Menü \menu{View}; geschlossen
öffnet es sich mit der nächsten neuen Tabelle wieder.

\subsection{Image}\label{sec:dock-image}
\index{Image-Fenster}

Das Fenster \gui{Image} zeigt Bilder aus den Rechnungen, etwa das Kamerabild des Raytracers
oder des Path Tracers. Es öffnet sich beim ersten Bild als schwebendes Fenster in der Mitte des
Hauptfensters und lässt sich wie jedes Dock andocken.

\begin{itemize}
  \item Achsen in physikalischen Einheiten, eine Farbskala am Rand.
  \item \gui{log}: logarithmische Farbskala, um schwaches Licht neben einem hellen Fleck zu
    sehen. Bei Farbbildern (z.\,B.\ Echtfarben-Bild mehrerer Wellenlängen) gibt es keine
    Skala; \gui{log} ist dann gesperrt.
  \item Unter dem Bild steht der Wert unter dem Mauszeiger (\enquote{x \dots{} y \dots{}} mit
    Wert und Einheit); ohne Mauszeiger eine Beschreibung des Bildes.
  \item \maus{Mausrad}: Zoom am Mauszeiger; \maus{Ziehen mit der linken Maustaste}:
    verschieben; \maus{Doppelklick}: wieder das ganze Bild.
  \item Gibt es mehrere Bilder, wählt eine Auswahlliste oben das angezeigte (\enquote{Lauf --
    Bild}). Ein neues Bild kommt automatisch nach vorn.
  \item \gui{Save…} (oben rechts): \gui{Image as PNG…}, \gui{View with axes as PNG…},
    \gui{Values as CSV…} (\autoref{sec:ergebnisse-speichern}).
\end{itemize}

\subsection{Plugin log}

Die Textausgabe aller Plugins, jede Zeile mit der Kennung des Plugins in eckigen Klammern,
z.\,B.\ \enquote{[raytracer] \dots}. Hier stehen Hinweise der Rechnungen und, wenn etwas schief
geht, die Fehlermeldung mit dem Python-Aufrufstapel. Das Dock liegt als Reiter neben
\gui{Plot} und öffnet sich von selbst, wenn eine Rechnung mit einem Fehler endet. Fehler, die
Sie selbst beheben können (falsche Eingabe, fehlende Auswahl), meldet das Plugin stattdessen
in einem Hinweisfenster. Das Protokoll behält die letzten 5000 Zeilen.

\subsection{Node workflow}

Das Fenster des Knoten-Workflows: links der Node-Editor mit Blöcken und Verbindungen, rechts
die Liste \gui{Branches} mit den Zweigen und darunter \gui{Steps}. Es erscheint, sobald ein
Dokument einen Workflow hat oder ein Plugin einen anlegt, und öffnet sich auch über
\menu{View,Node workflow}. Während einer Rechnung lautet der Titel \enquote{Node workflow --
running}. Anzeigeblöcke öffnen mit einem Doppelklick eigene Fenster \enquote{Display: \dots}
mit ihrem Bild oder Diagramm. Die Bedienung beschreibt \autoref{kap:workflow}.

\begin{hinweis}
Solange das Fenster \gui{Node workflow} den Tastaturfokus hat, wirken die Navigationstasten
(\keys{Z}, \keys{E} usw.) auf den Node-Editor und nicht auf die CAD-Ansichten.
\end{hinweis}

\section{Die Befehlszeile}\label{sec:befehlszeile}
\index{Befehlszeile}

Die Befehlszeile ist das Eingabefeld unter den Ansichten. Ohne laufenden Befehl zeigt sie
grau \enquote{Command:}. Sie nimmt vier Arten von Eingaben an, jeweils mit \keys{Enter}
abgeschlossen:

\begin{referenz}{>{\raggedright\arraybackslash}p{3.2cm}>{\raggedright\arraybackslash}p{3.0cm}X}
\toprule
\textbf{Eingabe} & \textbf{Beispiel} & \textbf{Wirkung} \\
\midrule
\endhead
Befehlsname oder Kurzform & \code{line}, \code{l}, \code{box} & startet den Befehl
  (\autoref{kap:befehle}). Unbekannte Namen meldet die Statusleiste: \enquote{Unknown command
  \dots}. \\
Punkt & \code{10,20,0}\newline \code{@5,0,0}\newline \code{25<30} & Koordinaten absolut, in der
  Konstruktionsebene, relativ zum letzten Punkt oder als Abstand und Winkel
  (\autoref{kap:modellieren}). \\
Zahl & \code{12.5} & Länge, Radius, Anzahl usw., wo der Befehl eine Zahl erwartet; sonst eine
  Länge in Richtung vom letzten Punkt zum Cursor. \\
leere Eingabe & \keys{Enter} & beendet den laufenden Befehl (z.\,B.\ eine Polylinie); ohne
  laufenden Befehl wird der letzte Befehl wiederholt. Ein \maus{Rechtsklick} ohne Ziehen in
  einer Ansicht wirkt ebenso. \\
\bottomrule
\end{referenz}

Während ein Befehl läuft, steht in der Befehlszeile die Aufforderung (\enquote{Prompt}), was
er als Nächstes erwartet, bei eingeschaltetem Ortho-Modus mit dem Zusatz
\enquote{[Ortho]}. Sie müssen die Befehlszeile nicht anklicken: Tippen Sie in einer Ansicht
Buchstaben oder Ziffern, wandern sie direkt in die Befehlszeile. \keys{Esc} bricht den Befehl
ab (\enquote{\dots{} cancelled}).

\begin{warnung}
Einige Tasten sind in den Ansichten mit Funktionen belegt und gelangen deshalb \emph{nicht} in
die Befehlszeile: \keys{Leertaste} maximiert die Ansicht, \keys{Z} leitet die Zoom-Kürzel ein.
Brauchen Sie diese Zeichen in der Befehlszeile, klicken Sie zuerst hinein. \keys{H} blendet die Auswahl aus, aber nur, wenn etwas ausgewählt ist; ohne Auswahl wird
\keys{H} wie jeder andere Buchstabe in die Befehlszeile geschrieben -- \code{h} und
\keys{Enter} starten dann die Helix.
\end{warnung}

Kann ein Befehl eine Eingabe nicht verwenden, nennt die Statusleiste den Grund, etwa
\enquote{Fillet: radius 5 too large for these curves}, und der Befehl wartet auf eine andere
Eingabe. Nur wenn der Befehl keinen eigenen Grund kennt, erscheint \enquote{Input not accepted
-- \textit{Aufforderung}} (\autoref{sec:mod-befehlszeile}).

Die Befehlszeile ist fest angedockt und lässt sich weder schließen noch verschieben.

\section{Die Statusleiste}\label{sec:statusleiste}
\index{Statusleiste}

Die Statusleiste am unteren Rand enthält links die Koordinaten und die aktive Ansicht, rechts
die Schalter für Ortho und Fang, während einer Rechnung deren Fortschritt und gegebenenfalls
den Hinweis auf den Entwicklermodus:

\begin{referenz}{>{\raggedright\arraybackslash}p{3.6cm}X}
\toprule
\textbf{Element} & \textbf{Bedeutung} \\
\midrule
\endhead
Koordinaten & \enquote{X \dots{} Y \dots{} Z \dots}: die Position des Mauszeigers in der
  Konstruktionsebene der Ansicht unter dem Zeiger, mit drei Nachkommastellen; rastet der
  Cursor auf einen Fangpunkt ein, dessen Position. Außerhalb der Ansichten Striche. \\
Aktive Ansicht & \enquote{Active view: \dots} \\
Meldungen & erscheinen links und überdecken dabei vorübergehend Koordinaten und aktive
  Ansicht: Meldungen der Befehle und Plugins (\enquote{Saved optik.rttf}, \enquote{12
  entities, 16928 triangles -- 0.91 s} nach dem Laden, Fortschrittstext einer Rechnung);
  beim Überfahren eines Objekts dessen Name, Art und ggf.\ Teilobjekt, z.\,B.\
  \enquote{part\_0\_1 $\cdot$ Shell $\cdot$ face 3}. \\
\gui{Ortho} & Ortho-Modus ein/aus (\keys{F8}): Punkte werden auf die Achsen der
  Konstruktionsebene durch den letzten Punkt gezogen. \\
\gui{Snap} & Objektfang insgesamt ein/aus (\keys{F3}). \\
\gui{Point}, \gui{End}, \gui{Corner}, \gui{Int}, \gui{Cen}, \gui{Quad}, \gui{Mid}, \gui{Edge},
  \gui{Near}, \gui{Grid} & die einzelnen Fangarten: Punkt, Endpunkt, Flächenecke,
  Schnittpunkt, Zentrum, Quadrantenpunkt, Mitte, Flächenkanten, nächster Kurvenpunkt,
  Rasterpunkt. Eingedrückt = aktiv. Anfangs sind alle außer \gui{Grid} aktiv. Einzelheiten in
  \autoref{kap:modellieren}. \\
Fortschrittsbalken und \gui{Cancel} & nur während einer Plugin-Rechnung sichtbar;
  \gui{Cancel} bricht sie ab (wie \keys{Esc}). Nach dem Ende meldet die Statusleiste
  \enquote{\textit{Plugin}: done} bzw.\ das Ergebnis, \enquote{\dots: cancelled} oder
  \enquote{\dots{} failed: \dots}. \\
\enquote{DEVELOPER MODE -- unsigned plugins allowed} & roter Balken ganz rechts, solange der
  Entwicklermodus eingeschaltet ist, in dem auch nicht signierte Plugins geladen werden
  (\autoref{kap:plugin-verwaltung}). \\
\bottomrule
\end{referenz}

Die Schalter \gui{Snap} und die Fangarten werden gespeichert.

\begin{warnung}
Der rote Balken \enquote{DEVELOPER MODE} bedeutet, dass Python-Code ohne geprüfte Herkunft
ausgeführt werden darf. Lassen Sie den Entwicklermodus nur eingeschaltet, wenn Sie selbst
Plugins entwickeln.
\end{warnung}

\section{Die Menüs}\label{sec:menues}
\index{Menü}

Die Menüleiste hat fünf Menüs: \menu{File}, \menu{Edit}, \menu{View}, \menu{Plugins} und
\menu{Help}. Auf dem Mac stehen \gui{Quit} und \gui{About Raytatouille Fusion} nicht in
\menu{File} bzw.\ \menu{Help}, sondern -- wie dort üblich -- im Programmmenü
\gui{Raytatouille Fusion}. Die Tastenkürzel in den folgenden Tabellen sind die unter Windows
und Linux; zum Mac siehe \autoref{sec:schreibweisen}.

\subsection{File}\label{sec:menu-file}

\begin{referenz}{>{\raggedright\arraybackslash}p{2.8cm}>{\raggedright\arraybackslash}p{2.4cm}X}
\toprule
\textbf{Eintrag} & \textbf{Kürzel} & \textbf{Wirkung} \\
\midrule
\endhead
\gui{New} & \keys{Strg+N} & Neues, leeres Dokument mit der Ebene \gui{Default}. Fragt vorher
  nach dem Speichern ungesicherter Änderungen. \\
\gui{Open...} & \keys{Strg+O} & Öffnet ein Dokument (\datei{.rttf}, \datei{.nfcad}) oder eine
  STEP-Datei (\datei{.step}, \datei{.stp}). Beides \emph{ersetzt} das aktuelle Dokument;
  vorher wird nach dem Speichern gefragt (\autoref{sec:oeffnen-speichern}). Ein Plugin-Paket
  (\datei{.nfplugin}) wird installiert. \\
\gui{Open Recent} & -- & Die letzten zehn geöffneten oder gespeicherten Dateien;
  \gui{Clear list} leert die Liste. \\
\gui{Import STEP...} & \keys{Strg+I} & Liest eine STEP-Datei und \emph{fügt} ihren Inhalt als
  Gruppe zum Dokument hinzu (\autoref{sec:step-import}). \\
\gui{Save} & \keys{Strg+S} & Speichert das Dokument unter seinem Namen; ein neues Dokument
  fragt nach dem Namen wie \gui{Save As...}. \\
\gui{Save As...} & \keys{Strg+Shift+S}\newline (unter Windows ohne Kürzel) & Speichert unter
  einem neuen Namen, immer als \datei{.rttf}. \\
\gui{Export STEP...} & -- & Schreibt die ausgewählten Objekte -- ohne Auswahl alle sichtbaren --
  als STEP (AP214); Punkte und Kurven werden nicht exportiert (\autoref{sec:step-export}). \\
\gui{Export STL...} & -- & Schreibt die Dreiecksnetze der ausgewählten bzw.\ aller sichtbaren
  Objekte als binäres STL (\autoref{sec:stl-export}). \\
\gui{Quit} & \keys{Strg+Q}\newline (unter Windows \keys{Alt+F4}) & Beendet das Programm; fragt
  vorher nach dem Speichern. Auf dem Mac im Programmmenü. \\
\bottomrule
\end{referenz}

\subsection{Edit}\label{sec:menu-edit}

\begin{referenz}{>{\raggedright\arraybackslash}p{3.3cm}>{\raggedright\arraybackslash}p{2.2cm}X}
\toprule
\textbf{Eintrag} & \textbf{Kürzel} & \textbf{Wirkung} \\
\midrule
\endhead
\gui{Undo} & \keys{Strg+Z} & Macht den letzten Schritt rückgängig. Der Eintrag nennt den
  Schritt, z.\,B.\ \enquote{Undo Delete 2 object(s)}; ist nichts rückgängig zu machen, ist er
  gesperrt (\autoref{sec:undo}). \\
\gui{Redo} & \keys{Strg+Shift+Z}\newline (Windows auch \keys{Strg+Y}) & Stellt den zuletzt
  rückgängig gemachten Schritt wieder her. \\
\gui{Delete} & \keys{Entf}, \keys{Rück} & Löscht die ausgewählten Objekte (ein
  Rückgängig-Schritt). Flächen, Kanten und Eckpunkte allein lassen sich nicht löschen
  (\enquote{Faces, edges and vertices cannot be deleted on their own}). \\
\gui{Select All} & \keys{Strg+A} & Wählt alle sichtbaren Objekte aus. \\
\gui{Group...} & \keys{Strg+G} & Fasst die ausgewählten Objekte zu einer neuen Gruppe
  zusammen; fragt nach dem Namen (\autoref{sec:gruppen}). \\
\gui{Ungroup} & \keys{Strg+Shift+G} & Löst die im Modellbaum markierte(n) Gruppe(n) auf. \\
\gui{Hide} & \keys{H} & Blendet die ausgewählten Objekte aus und hebt ihre Auswahl auf
  (\autoref{sec:ausblenden}). \\
\gui{Isolate selected} & -- & Blendet alle Objekte außer den ausgewählten aus. \\
\gui{Unhide} & \keys{Shift+H} & Zeigt alle mit \gui{Hide} oder \gui{Isolate selected}
  ausgeblendeten Objekte wieder; Ebenen und Gruppen behalten ihre Schalter. \\
\gui{Show all} & -- & Macht alles sichtbar: Objekte, Gruppen und Ebenen. \\
\gui{Select None / Cancel} & \keys{Esc} & Bricht den laufenden Befehl ab; läuft keiner, die
  laufende Plugin-Rechnung; sonst hebt es die Auswahl auf. \\
\bottomrule
\end{referenz}

\subsection{View}\label{sec:menu-view}

\begin{referenz}{>{\raggedright\arraybackslash}p{3.0cm}>{\raggedright\arraybackslash}p{2.2cm}X}
\toprule
\textbf{Eintrag} & \textbf{Kürzel} & \textbf{Wirkung} \\
\midrule
\endhead
\gui{Wireframe}, \gui{Surfaces}, \gui{Shaded} & -- & Darstellungsmodus der aktiven Ansicht
  (\autoref{sec:darstellungsmodi}); der aktuelle ist abgehakt. Mit \keys{Shift} gewählt: alle
  Ansichten. \\
\gui{Zoom extents} & \keys{Z}, \keys{E} & ganzes Modell in der aktiven Ansicht. \\
\gui{Zoom selected} & \keys{Z}, \keys{S} & Auswahl füllt die aktive Ansicht. \\
\gui{Tools} & -- & Werkzeugspalte ein/aus. \\
\gui{Model} & -- & Dock Modellbaum ein/aus. \\
\gui{Results} & -- & Dock Ergebnisse ein/aus. \\
\gui{Node workflow} & -- & Fenster des Knoten-Workflows ein/aus. \\
\gui{Image} & -- & Bildfenster ein/aus. \\
\gui{Plot} & -- & Diagramm ein/aus. \\
\gui{Properties} & -- & Dock Eigenschaften ein/aus. \\
\gui{Layers} & -- & Dock Ebenen ein/aus. \\
\gui{Plugin log} & -- & Protokoll der Plugins ein/aus. \\
\gui{Plugin} & -- & Bedienseite der Plugins ein/aus. \\
\bottomrule
\end{referenz}

Die Dock-Einträge sind abgehakt, solange das Dock geöffnet ist. Ein Dock, das als Reiter
hinter einem anderen liegt, gilt als geöffnet; klicken Sie dann auf seinen Reiter.

\subsection{Plugins}\label{sec:menu-plugins}

\begin{referenz}{>{\raggedright\arraybackslash}p{3.6cm}X}
\toprule
\textbf{Eintrag} & \textbf{Wirkung} \\
\midrule
\endhead
\gui{Settings...} & Dialog \gui{Plugin settings}: Python-Interpreter, Umgebung prüfen und
  anlegen, Plugin-Ordner, Entwicklermodus (\autoref{sec:python-umgebung},
  \autoref{kap:plugin-verwaltung}). \\
\gui{Install plugin...} & Installiert ein Plugin-Paket (\datei{.nfplugin}). Danach meldet ein
  Fenster das Ergebnis und gegebenenfalls fehlende Python-Pakete mit dem
  \code{pip}-Befehl. \\
\gui{Uninstall plugin} $\triangleright$ & Untermenü mit allen Plugins (Name und Version).
  Mit \gui{Install plugin...} installierte lassen sich hier nach einer Rückfrage entfernen.
  Plugins aus einem Plugin-Ordner -- auch die mitgelieferten -- sind gesperrt (\enquote{\dots{}
  -- from a plugin folder}); der Tooltip nennt den Ordner. \\
\gui{Reload plugins} & Beendet alle Plugin-Prozesse, liest die Plugin-Ordner neu ein und
  startet die Plugins neu -- etwa nach einem Update eines Rechenkerns. \\
\gui{Show in the side panels} $\triangleright$ & Ein Häkchen je Plugin: ob seine Gruppe in der
  Werkzeugspalte und seine Seite im Dock \gui{Plugin} erscheinen. Die Aktionen bleiben im Menü
  \menu{Plugins}. Die Wahl wird gespeichert. \\
\textit{je Plugin ein Untermenü} $\triangleright$ & Die Aktionen des Plugins, z.\,B.\ unter
  \gui{Raytracer} \enquote{Load system}, \enquote{Trace rays} \dots{} Der Tooltip des
  Untermenüs zeigt die Beschreibung des Plugins. Startet das Plugin gerade, steht dort
  \enquote{(starting \dots)}, läuft es nicht, \enquote{(not running -- see the plugin log)}. \\
\textit{(no plugins loaded -- see Settings)} & erscheint, wenn kein einziges Plugin geladen
  ist. \\
\bottomrule
\end{referenz}

Aktionen, die zu einem bestimmten Block des Workflows gehören, stehen nicht hier, sondern im
Kontextmenü dieses Blocks im Node-Editor (\autoref{kap:workflow}).

\begin{hinweis}
Es kann immer nur eine Plugin-Rechnung gleichzeitig laufen. Startet man eine zweite, meldet
die Statusleiste \enquote{A plugin job is still running}. Ausnahme: Plugins, die dem Workflow
folgen, zeichnen kurz nach jeder Änderung im Node-Editor neu, sofern sie gerade nichts zu tun
haben.
\end{hinweis}

\subsection{Help}\label{sec:menu-help}

\begin{referenz}{>{\raggedright\arraybackslash}p{4.6cm}X}
\toprule
\textbf{Eintrag} & \textbf{Wirkung} \\
\midrule
\endhead
\gui{Demo results on the selected object} & Erzeugt zu Testzwecken einen Lauf \enquote{Demo
  run} mit Beispielergebnissen auf dem ersten ausgewählten Körper bzw.\ der ersten
  ausgewählten Fläche: Strahlen (\enquote{Rays}) von einer Quelle über der ersten Fläche, ein
  Feld darauf (\enquote{Irradiance}) und eine Kurve (\enquote{Merit}). Nützlich,
  um die Ergebnisanzeige ohne Python-Umgebung auszuprobieren. Ohne geeignete Auswahl meldet
  die Statusleiste \enquote{Select a surface or solid first}. \\
\gui{Check for updates...} & Fragt nach neuen Versionen des Programms und der Plugins
  (\autoref{sec:updates}). \\
\gui{About Raytatouille Fusion} & Programmversion, Version und Selbsttest des NURBS-Kerns
  Nurbforge, Qt-Version. Auf dem Mac im Programmmenü. \\
\bottomrule
\end{referenz}

\subsection{Kontextmenüs}

Neben der Menüleiste gibt es Kontextmenüs (Rechtsklick) an diesen Stellen:

\begin{itemize}
  \item \textbf{Kopfzeile einer Ansicht} (Linksklick): Standardansicht, Projektion, Zoom,
    Maximieren (\autoref{sec:ansichtsmenue}).
  \item \textbf{Modellbaum}: \gui{Group selection...}, \gui{Ungroup}, \gui{Hide},
    \gui{Isolate}, \gui{Unhide}, \gui{Show all} (\autoref{sec:modellbaum}).
  \item \textbf{Dock Results}: \gui{Delete run}, \gui{Convert to CAD curves...}
    (\autoref{sec:dock-results}).
  \item \textbf{Node-Editor}: Blöcke, Verbindungen, Gruppen und Varianten
    (\autoref{kap:workflow}).
\end{itemize}

Ein Rechtsklick in eine Ansicht öffnet kein Menü: Mit der rechten Maustaste wird gedreht
oder verschoben. Ein kurzer Rechtsklick ohne Ziehen wirkt wie \keys{Enter}: Während ein Befehl
läuft, beendet er die Eingabe; läuft keiner, wiederholt er den letzten Befehl.
